% area: Category Theory

\begin{definition}[Category] \label{def:category} 
A \emph{category} $\mathcal{C}$ consists of the following data:
\begin{enumerate} 
  \item A collection of objects, denoted $\operatorname{Ob}(\mathcal{C})$. 
  \item For every pair of objects $X,Y \in \operatorname{Ob}(\mathcal{C})$, a set $$ \operatorname{Hom}_{\mathcal{C}}(X,Y) $$ whose elements are called \emph{morphisms} (or arrows) from $X$ to $Y$. 
  \item For every triple of objects $X,Y,Z$, a composition map $$ \operatorname{Hom}_{\mathcal{C}}(Y,Z) \times \operatorname{Hom}_{\mathcal{C}}(X,Y) \to \operatorname{Hom}_{\mathcal{C}}(X,Z), $$ written $$ (g,f) \mapsto g \circ f . $$ 
  \item For every object $X$, an identity morphism $$ \operatorname{id}_X \in \operatorname{Hom}_{\mathcal{C}}(X,X). $$ 
  \end{enumerate} 
These data satisfy the following axioms: 
\begin{enumerate} 
\item \textbf{Associativity:} For morphisms $$ X \xrightarrow{f} Y \xrightarrow{g} Z \xrightarrow{h} W $$ we have $$ h \circ (g \circ f) = (h \circ g) \circ f . $$ 
\item \textbf{Identity:} For every morphism $f : X \to Y$, $$ f \circ \operatorname{id}_X = f \qquad\text{and}\qquad \operatorname{id}_Y \circ f = f . $$ \end{enumerate}
\end{definition}



\begin{remark}[Notation for Morphisms]
\label{rmk:morphism-notation}

Let $\mathcal{C}$ be a category as in Definition~\ref{def:category}.  
If $f \in \Hom_{\mathcal{C}}(X,Y)$ we write
$$
f : X \to Y.
$$

The object $X$ is called the \emph{domain} of $f$ and $Y$ the \emph{codomain}.

Composition of morphisms is written
$$
g \circ f : X \to Z
$$
whenever $f : X \to Y$ and $g : Y \to Z$ as in Definition~\ref{def:category}.
\end{remark}


\begin{definition}[Isomorphism]
\label{def:isomorphism}

Let $\mathcal{C}$ be a category (Definition~\ref{def:category}).  

A morphism
$$
f : X \to Y
$$
is called an \emph{isomorphism} if there exists a morphism
$$
g : Y \to X
$$
such that
$$
g \circ f = \operatorname{id}_X
\qquad
f \circ g = \operatorname{id}_Y.
$$

The morphism $g$ is called the \emph{inverse} of $f$ and is denoted
$$
f^{-1}.
$$

If such a morphism exists we say that $X$ and $Y$ are \emph{isomorphic} and write
$$
X \cong Y.
$$

This notion uses the identity morphisms and composition introduced in
Definition~\ref{def:category}.
\end{definition}


\begin{proposition}[Uniqueness of Inverses]
\label{prop:inverse-unique}

Let $\mathcal{C}$ be a category (Definition~\ref{def:category}).  
If a morphism $f : X \to Y$ is an isomorphism (Definition~\ref{def:isomorphism}),
then its inverse is unique.
\end{proposition}

\begin{proof}

Suppose
$$
g,h : Y \to X
$$
are both inverses of $f$. Then

$$
g \circ f = \operatorname{id}_X,
\qquad
f \circ g = \operatorname{id}_Y,
$$

and

$$
h \circ f = \operatorname{id}_X,
\qquad
f \circ h = \operatorname{id}_Y.
$$

Using associativity of composition from Definition~\ref{def:category} we obtain

$$
g
=
g \circ \operatorname{id}_Y
=
g \circ (f \circ h)
=
(g \circ f) \circ h
=
\operatorname{id}_X \circ h
=
h.
$$

Thus the inverse is unique.
\end{proof}


\begin{corollary}
\label{cor:iso-equivalence-relation}

Let $\mathcal{C}$ be a category (Definition~\ref{def:category}).  
The relation
$$
X \cong Y
$$
defined via isomorphisms (Definition~\ref{def:isomorphism})
is an equivalence relation on $\operatorname{Ob}(\mathcal{C})$.
\end{corollary}

\begin{proof}

\textbf{Reflexivity.}  
For every object $X$, the identity morphism $\operatorname{id}_X$ from
Definition~\ref{def:category} satisfies

$$
\operatorname{id}_X \circ \operatorname{id}_X = \operatorname{id}_X.
$$

Thus $\operatorname{id}_X$ is an isomorphism.

\textbf{Symmetry.}  
If $f : X \to Y$ is an isomorphism, Definition~\ref{def:isomorphism}
provides an inverse $f^{-1} : Y \to X$.

\textbf{Transitivity.}  
If $f : X \to Y$ and $g : Y \to Z$ are isomorphisms, then

$$
(g \circ f)^{-1} = f^{-1} \circ g^{-1}
$$

by associativity of composition in Definition~\ref{def:category}.
\end{proof}


\begin{definition}[Subcategory]
\label{def:subcategory}

Let $\mathcal{C}$ be a category (Definition~\ref{def:category}).  
A \emph{subcategory} $\mathcal{D}$ of $\mathcal{C}$ consists of

\begin{enumerate}

\item a collection of objects
$$
\operatorname{Ob}(\mathcal{D}) \subseteq \operatorname{Ob}(\mathcal{C})
$$

\item for each pair of objects $X,Y \in \operatorname{Ob}(\mathcal{D})$,
a subset

$$
\Hom_{\mathcal{D}}(X,Y)
\subseteq
\Hom_{\mathcal{C}}(X,Y)
$$

\end{enumerate}

such that

\begin{enumerate}

\item identity morphisms from Definition~\ref{def:category} belong to
$\mathcal{D}$,

\item composition in $\mathcal{D}$ is the restriction of composition in
$\mathcal{C}$.

\end{enumerate}

\end{definition}


\begin{definition}[Opposite Category]
\label{def:opposite-category}

Let $\mathcal{C}$ be a category (Definition~\ref{def:category}).  
The \emph{opposite category} $\mathcal{C}^{op}$ is defined as follows.

\begin{enumerate}

\item Objects:
$$
\operatorname{Ob}(\mathcal{C}^{op}) = \operatorname{Ob}(\mathcal{C})
$$

\item Morphisms:
$$
\Hom_{\mathcal{C}^{op}}(X,Y)
=
\Hom_{\mathcal{C}}(Y,X)
$$

\item Composition is defined by reversing composition in $\mathcal{C}$.

\end{enumerate}

Associativity and identity laws follow directly from those in
Definition~\ref{def:category}.
\end{definition}


\begin{remark}[Duality Principle]
\label{rmk:duality}

Let $\mathcal{C}$ be a category (Definition~\ref{def:category}).  
Any statement about $\mathcal{C}$ has a \emph{dual statement}
obtained by replacing

\begin{itemize}

\item morphisms $f : X \to Y$ by $f : Y \to X$

\item compositions $g \circ f$ by $f \circ g$

\end{itemize}

which corresponds precisely to passing to the opposite category
$\mathcal{C}^{op}$ defined in Definition~\ref{def:opposite-category}.

This observation is called the \emph{principle of duality}.
\end{remark}


\begin{definition}[Functor]
\label{def:functor}
  A functor $F$ is a map of categories that preserves commutative diagrams. In particular, given two categories, \ref{def:category} $\mathcal{C}$, $\mathcal{D}$, a functor 
 $F : \mathcal{C} \to \mathcal{D}$ satisfies the following properties:

\begin{enumerate}
  \item Associates each object $X \in \operatorname{Ob}(\mathcal{C})$ to an object $F(X) \in \operatorname{Ob}(\mathcal{D})$ 
  \item Associates each morphism $f : X \to Y$ in $\mathcal{C}$ in $\mathcal{C}$ to a morphism $F(f) : F(X) \to F(Y)$ in $\mathcal{D}$ such that the following conditions hold: 
    \begin{enumerate}
      \item $F(\operatorname{id}_X) = \operatorname{id}_{F(X)}$
      \item $F(g \circ f) = F(g) \circ F(f)$ for all morphisms $f : X \to Y$ and $g : Y \to Z$ in $\mathcal{C}$.  
    \end{enumerate}
\end{enumerate}
That is, functors must preserve identity morphisms and composition of morphisms.
\end{definition}

\begin{lemma}[Functors Preserve Isomorphisms]
\label{lem:functor-preserve-iso}

Let 
$$
F : \mathcal{C} \to \mathcal{D}
$$
be a functor (Definition~\ref{def:functor}).  

If
$$
f : X \to Y
$$
is an isomorphism in $\mathcal{C}$ (Definition~\ref{def:isomorphism}),
then
$$
F(f) : F(X) \to F(Y)
$$
is an isomorphism in $\mathcal{D}$.

References: Definition~\ref{def:functor}, Definition~\ref{def:isomorphism}.
\end{lemma}

\begin{proof}

Since $f$ is an isomorphism, there exists a morphism
$$
f^{-1} : Y \to X
$$
such that
$$
f^{-1} \circ f = \operatorname{id}_X,
\qquad
f \circ f^{-1} = \operatorname{id}_Y .
$$

Apply the functor $F$ to the first equality:

$$
F(f^{-1} \circ f) = F(\operatorname{id}_X).
$$

By functoriality of composition and identities this becomes

$$
F(f^{-1}) \circ F(f) = \operatorname{id}_{F(X)}.
$$

Applying $F$ to the second equality gives

$$
F(f) \circ F(f^{-1}) = \operatorname{id}_{F(Y)}.
$$

Thus $F(f^{-1})$ is an inverse of $F(f)$, and therefore $F(f)$ is an isomorphism.
\end{proof}

\begin{definition}
\label{def:functor-composition}

Let
$$
F : \mathcal{C} \to \mathcal{D},
\qquad
G : \mathcal{D} \to \mathcal{E}
$$
be functors (Definition~\ref{def:functor}).

The \emph{composition of functors}
$$
G \circ F : \mathcal{C} \to \mathcal{E}
$$
is defined as follows.

\textbf{On objects}
$$
(G \circ F)(X) = G(F(X)).
$$

\textbf{On morphisms}

For
$$
f : X \to Y
$$

define

$$
(G \circ F)(f) = G(F(f)).
$$

\end{definition}

\begin{lemma}
\label{lem:functor-composition}

The composition of functors is a functor.

References: Definition~\ref{def:functor}, Definition~\ref{def:functor-composition}.
\end{lemma}

\begin{proof}

Let

$$
F : \mathcal{C} \to \mathcal{D},
\qquad
G : \mathcal{D} \to \mathcal{E}.
$$

First we check preservation of identities.

For an object $X$,

$$
(G \circ F)(\operatorname{id}_X)
=
G(F(\operatorname{id}_X)).
$$

Since $F$ preserves identities,

$$
F(\operatorname{id}_X) = \operatorname{id}_{F(X)}.
$$

Thus

$$
(G \circ F)(\operatorname{id}_X)
=
G(\operatorname{id}_{F(X)})
=
\operatorname{id}_{G(F(X))}.
$$

Next we check preservation of composition.

Let

$$
f : X \to Y,
\qquad
g : Y \to Z.
$$

Then

$$
(G \circ F)(g \circ f)
=
G(F(g \circ f)).
$$

Since $F$ preserves composition,

$$
F(g \circ f) = F(g) \circ F(f).
$$

Applying $G$ gives

$$
G(F(g) \circ F(f))
=
G(F(g)) \circ G(F(f)).
$$

Thus

$$
(G \circ F)(g \circ f)
=
(G \circ F)(g) \circ (G \circ F)(f).
$$

Therefore $G \circ F$ is a functor.

\end{proof}

\begin{definition}
\label{def:identity-functor}

Let $\mathcal{C}$ be a category.

The \emph{identity functor}

$$
\operatorname{Id}_{\mathcal{C}} : \mathcal{C} \to \mathcal{C}
$$

is defined by

\textbf{Objects}

$$
\operatorname{Id}_{\mathcal{C}}(X) = X
$$

\textbf{Morphisms}

$$
\operatorname{Id}_{\mathcal{C}}(f) = f.
$$

\end{definition}

\begin{lemma}
\label{lem:identity-functor}

The identity functor is a functor.

References: Definition~\ref{def:functor}, Definition~\ref{def:identity-functor}.
\end{lemma}

\begin{proof}

Let $f : X \to Y$ and $g : Y \to Z$.

Then

$$
\operatorname{Id}_{\mathcal{C}}(g \circ f)
=
g \circ f
=
\operatorname{Id}_{\mathcal{C}}(g) \circ \operatorname{Id}_{\mathcal{C}}(f).
$$

Also

$$
\operatorname{Id}_{\mathcal{C}}(\operatorname{id}_X)
=
\operatorname{id}_X.
$$

Thus the functor axioms hold.

\end{proof}


\begin{proposition}
\label{prop:functor-associative}

Functor composition is associative.

References: Definition~\ref{def:functor-composition}.
\end{proposition}

\begin{proof}

Let

$$
F : \mathcal{C} \to \mathcal{D},
\qquad
G : \mathcal{D} \to \mathcal{E},
\qquad
H : \mathcal{E} \to \mathcal{F}.
$$

For an object $X$,

$$
H \circ (G \circ F)(X)
=
H(G(F(X))).
$$

Also

$$
(H \circ G) \circ F (X)
=
H(G(F(X))).
$$

Thus the two functors agree on objects.

For a morphism $f : X \to Y$,

$$
H \circ (G \circ F)(f)
=
H(G(F(f))).
$$

Similarly,

$$
(H \circ G) \circ F (f)
=
H(G(F(f))).
$$

Thus the functors agree on morphisms.

\end{proof}


\begin{example}
\label{ex:forgetful}

Let

$$
\mathbf{Grp}
$$

be the category of groups and group homomorphisms, and

$$
\mathbf{Set}
$$

the category of sets.

Define a functor

$$
U : \mathbf{Grp} \to \mathbf{Set}
$$

as follows.

\textbf{Objects}

$$
U(G) = \text{the underlying set of } G.
$$

\textbf{Morphisms}

For a group homomorphism

$$
f : G \to H
$$

define

$$
U(f) = f
$$

viewed as a function between sets.

This functor is called the \emph{forgetful functor}.

\end{example}

\begin{definition}
\label{def:natural-transformation}

Let

$$
F,G : \mathcal{C} \to \mathcal{D}
$$

be functors (Definition~\ref{def:functor}).

A \emph{natural transformation}

$$
\eta : F \Rightarrow G
$$

consists of morphisms

$$
\eta_X : F(X) \to G(X)
$$

for every object $X$ of $\mathcal{C}$ such that for every morphism

$$
f : X \to Y
$$

the following diagram commutes:

$$
\begin{tikzcd}
F(X) \arrow[r,"\eta_X"] \arrow[d,"F(f)"']
&
G(X) \arrow[d,"G(f)"]
\\
F(Y) \arrow[r,"\eta_Y"]
&
G(Y)
\end{tikzcd}
$$

\end{definition}
